% Part: computability
% Chapter: computability-theory
% Section: complete-ce-sets

\documentclass[../../../include/open-logic-section]{subfiles}

\begin{document}

\olfileid{cmp}{thy}{cce}
\olsection{بشپړ په محاسبوي ډول د شمېر وړ سټونه}

\begin{defn}
سټ~$A$ (د ډېر-پر-يو راکمېدنې لاندې) يو
\emph{بشپړ !!{computably enumerable} سټ} دے، که
\begin{enumerate}
\item $A$ په محاسبوي ډول د شمېر وړ وي، او
\item د هر بل په محاسبوي ډول د شمېر وړ سټ~$B$ دپاره $B \leq_m A$ وي۔
\end{enumerate}
\end{defn}

په بله وينا، بشپړ په محاسبوي ډول د شمېر وړ سټونه د ټولو ممکنه په
محاسبوي ډول د شمېر وړ سټونو ``تر ټولو ستونزمن'' دي۔ دوي اجازه راکوي
چې د \emph{هر} په محاسبوي ډول د شمېر وړ سټ په اړه پوښتنو ته ځواب
ووايو۔

\begin{thm}
$K$، $K_0$ او~$K_1$ ټول بشپړ په محاسبوي ډول د شمېر وړ سټونه دي۔
\end{thm}

\begin{proof}
د دې د ليدلو دپاره چې~$K_0$ بشپړ دے، $B$ هر په محاسبوي ډول د شمېر
وړ سټ وي۔ نو د کوم شاخص~$e$ دپاره
\[
B = W_e = \Setabs{x}{\cfind{e}(x) \fdefined}.
\]
$f$ هغه تابع وي چې $f(x) = \tuple{e, x}$۔ بيا د هر طبيعي عدد~$x$
دپاره $x \in B$ هله او يوازې هله چې $f(x) \in K_0$ وي۔ په بله
وينا، $f$، $B$، $K_0$ ته راکموي۔

د دې د ليدلو دپاره چې~$K_1$ بشپړ دے، پام وکړئ چې د
\olref[k1]{prop:k1} په ثبوت کښې مو~$K_0$ ورته راکم کړے و۔ نو د
\olref[ppr]{prop:trans-red} له مخې هر په محاسبوي ډول د شمېر وړ سټ
~$K_1$ ته هم راکمول کېدے شي۔

$K_0$ هم په ډېره ورته طريقه~$K$ ته راکمول کېدے شي۔
\textbf{د ژباړې د نسخې سمون (OLCMP-034):} د سرچينې وروستۍ جمله د
راکمونې لورے سرچپه ليکلے و: له قطري سټ څخه د جوړو سټ ته راکمونه د
لومړي سټ بشپړتيا نۀ ثابتوي۔ دلته د جوړو بشپړ سټ قطري سټ ته راکمېږي،
څو د انتقال له مخې هر c.e. سټ هم قطري سټ ته راکم شي۔
\end{proof}

\begin{prob}
له~$K$ څخه~$K_0$ ته يوه راکمونه ورکړئ۔
\end{prob}

\begin{digress}
نو څرګنده شوه چې تر اوسه مو د په محاسبوي ډول د شمېر وړ سټونو هره
کتلې بېلګه يا محاسبه کېدونکې ده، يا بشپړه۔ دا بايد عجيبه وبرېښي!{}
ايا داسې په محاسبوي ډول د شمېر وړ سټونه شته چې نۀ محاسبه کېدونکي او
نۀ بشپړ وي؟ ځواب هو دے، خو دا تر هغه وخته نۀ وه ثابته شوې چې د
۱۹۵۰مو کلونو په منځ کښې فريډبرګ او موچنيک په خپلواکه توګه ثابته کړه۔
\end{digress}

\end{document}
